\chapter{Discussion}

\section{Reddening Comparison}
\cite{2009ApJ...696.1498M}


\section{Distance Comparison}

\section{Luminosity Comparison}

\begin{figure}[h]
	\centering
	\includegraphics[width=\linewidth]{M_V}
\end{figure}

\subsubsection*{Effect of Metallicity}
While metallicity can affect the zero-point and slope of the period--luminosity relation, it is not explicitly considered in this thesis. However, the framework developed here is readily extendable to incorporate metallicity corrections in future studies, enabling a more comprehensive calibration of the Leavitt Law across different galactic environments.

\subsubsection*{Target-Based Reddening Law}
The refined calibration methodology also provides insights into the reddening law for individual Cepheids. By comparing observed deviations across multiple bands and Wesenheit functions, star-specific reddening ratios can be inferred, allowing a more precise correction of extinction effects. This target-based approach can improve the accuracy of distance estimates, particularly in regions with non-standard or spatially varying extinction, and serves as a foundation for further refinement of both Galactic and extragalactic distance measurements.


\cite{2009MNRAS.396.1287S}